% Gerado por scripts/etapa1_sazonalidade_anual.py — requer \usepackage{booktabs}
\begin{table}[htbp]
\centering
\small
\caption{Eventos de calendário com efeito consistente sobre a demanda}
\label{tab:saz_eventos_resumo}
\begin{tabular}{lrrr}
\toprule
Evento & Tipo & Efeito ajustado & Ocorrências no mesmo sentido \\
\midrule
Dia do Trabalho & Nacional & +17,8\% & 5 de 5 \\
Dia dos Presidentes & Nacional & +12,3\% & 6 de 6 \\
Memorial Day & Nacional & +9,1\% & 5 de 5 \\
Martin Luther King & Nacional & +7,4\% & 5 de 5 \\
Dia dos Veteranos & Nacional & +5,9\% & 4 de 5 \\
Páscoa Ortodoxa & Religioso & -5,3\% & 4 de 5 \\
Dia dos Pais & Cultural & -5,5\% & 4 de 5 \\
Halloween & Cultural & -6,5\% & 4 de 5 \\
Dia dos Namorados & Cultural & -7,0\% & 6 de 6 \\
Páscoa & Cultural & -8,8\% & 5 de 6 \\
Cinco de Mayo & Cultural & -9,6\% & 5 de 5 \\
Dia das Mães & Cultural & -13,3\% & 5 de 5 \\
Ano Novo & Nacional & -25,9\% & 5 de 5 \\
Ação de Graças & Nacional & -26,1\% & 5 de 5 \\
\bottomrule
\end{tabular}
\par\smallskip
{\footnotesize Nota: FOODS, lojas da Califórnia, dias disponíveis e de loja aberta; demanda média por série-dia. Efeito ajustado: diferença \% entre a demanda no dia do evento e a média dos dias de mesmo dia da semana e mesma condição de SNAP, numa janela de até 6 semanas antes ou depois do evento (excluídos os 3 dias anteriores e posteriores); média das ocorrências de 2011 a 2016. Incluídos os eventos com efeito de pelo menos 5\% em módulo que se repete no mesmo sentido em pelo menos 4 ocorrências. Tabela completa no apêndice.}
\end{table}
